% Draft: methods paper on personalised MANO fits from sparse Vicon hand markers during hand-object interaction.
% Rule for this draft: only numbers supplied by the authors appear; everything else is \tbd{}.
% Build: latexmk -pdf main.tex
\documentclass[11pt,a4paper]{article}
\usepackage[margin=2.5cm]{geometry}
\usepackage[T1]{fontenc}
\usepackage{lmodern}
\usepackage{amsmath,amssymb}
\usepackage{booktabs,tabularx,multirow}
\usepackage{graphicx}
\usepackage{xcolor}
\usepackage{enumitem}
\usepackage{listings}
\usepackage[round,authoryear]{natbib}
\usepackage[colorlinks,linkcolor=blue!50!black,citecolor=blue!50!black,urlcolor=blue!50!black]{hyperref}

\newcommand{\tbd}[1]{\textcolor{red}{\textbf{[TBD: #1]}}}
\newcommand{\note}[1]{\textcolor{orange!80!black}{\textit{[Note: #1]}}}
\lstset{basicstyle=\ttfamily\small,breaklines=true,columns=fullflexible,frame=single,framesep=4pt}

\title{Personalised MANO hand reconstruction from sparse optical markers during hand--object interaction:\\
a validated, reproducible pipeline for labelling, gap filling and fitting}
\author{\tbd{First author}, \tbd{co-authors}\\ \small \tbd{affiliations}}
\date{Draft, \today}

\begin{document}
\maketitle

% ---------------------------------------------------------------------------------------------------------------
\section*{Title options}
\begin{enumerate}[nosep]
  \item Personalised MANO hand reconstruction from sparse optical markers during hand--object interaction: a validated, reproducible pipeline (current).
  \item From mislabelled markers to joint angles: learned labelling, gap filling and personalised hand-model fitting for sparse Vicon hand capture.
  \item Where does sparse-marker hand capture fail during object interaction? A stage-by-stage validated pipeline with participant-specific learned labelling and filling.
  \item Sparse-marker hand motion capture during grasping: participant-specific learned labelling and gap filling for personalised MANO fits.
\end{enumerate}

% ---------------------------------------------------------------------------------------------------------------
\begin{abstract}
\noindent
Optical hand motion capture with a sparse clinical marker set (22 markers, 200\,Hz) fails in three ways during
hand--object interaction (HOI): automatic labels are unreliable, grasps cause long occlusions, and joint angles
require a personalised hand model. We present and validate a pipeline that addresses each. A MANO hand model is fitted to
each participant's static trial (shape, scale, pose and per-marker surface attachments) and reviewed by hand. A
SOMA-style transformer labeller, trained on synthetic clouds from a participant-calibrated library, relabels
every frame. A multilayer perceptron, trained on the participant's own labelled frames by hiding observed markers,
fills gaps in a hand-aligned frame. A per-frame pose fit then holds the personalised shape fixed and down-weights
filled values. In 17 recorded participants, only 0.1--15\% of frames per trial were geometrically
self-consistent under the original labels, and HOI trials were 2--10$\times$ worse than hands-only trials.
On artificial gaps in a held-out HOI trial, the participant-specific filler had a lower worst-case error than
rigid-body filling for seven of eight participants (e.g.\ 44.9 to 9.8\,mm for 10\,s gaps). Across 34 trials of nine
participants it left a median 10.2\% of missing values unfilled, against 47.6\% for rigid-body filling. With filled
values weighted 0.1, the fit error on observed markers was never worse than with rigid-body filling in hands-only
trials (largest per-frame error, median change $-0.4$\,mm over 17 trials), but was worse in 7 of 17 HOI trials, by up
to 0.7\,mm. Labelling accuracy \tbd{SOMA vs Vicon on 500 verified frames} and gap-fill accuracy inside verified grasp
windows \tbd{} complete the validation. Training is seeded, results carry provenance records, and code is open.
\end{abstract}
\note{About 250 words; tighten once the TBD numbers are in. Cohort numbers from Tables~\ref{tab:gapfill-cohort} and~\ref{tab:fit-cohort} and the Completeness section (measured 2026-10-29).}

% ---------------------------------------------------------------------------------------------------------------
\section{Introduction}
\begin{itemize}[nosep]
  \item \textbf{Need.} Finger kinematics during object manipulation matter for rehabilitation, ergonomics and motor control; optical capture with a sparse marker set is the practical laboratory standard.
  \item \textbf{Failure point 1: labels.} Vicon's automatic labelling breaks during movement. In our 17-participant data set, only 0.1--15\% of frames per trial are fully geometrically self-consistent, and HOI trials are 2--10$\times$ worse than hands-only trials.
  \item \textbf{Failure point 2: gaps.} Grasps occlude markers for long periods; standard pattern and rigid-body filling degrade on long gaps and leave many values empty.
  \item \textbf{Failure point 3: personalisation.} Sparse markers under-determine joint angles; a hand model with participant-specific shape and marker attachments is needed.
  \item \textbf{Contribution.} Not a new network, but a validated, reproducible combination for sparse hand markers during HOI: (i) a hand-reviewed personalised static MANO fit, (ii) participant-calibrated learned labelling, (iii) participant-specific learned gap filling, (iv) a movement fit that down-weights filled values, (v) stage-by-stage validation, including where a stage does \emph{not} help, and (vi) a tutorial with provenance records.
\end{itemize}

% ---------------------------------------------------------------------------------------------------------------
\section{Related work}
\begin{itemize}[nosep]
  \item \textbf{Hand models.} MANO \citep{romero2017mano}: shape and pose space learned from dense scans.
  \item \textbf{Model fitting to markers.} MoSh \citep{loper2014mosh} and MoSh++/AMASS \citep{mahmood2019amass} fit body (and SMPL+H hand) models with marker attachments; GRAB \citep{taheri2020grab} captures whole-body grasps with dense hand markers. Our setting differs: 3 markers per finger, a clinical protocol, per-participant review.
  \item \textbf{Learned labelling.} SOMA \citep{ghorbani2021soma}: self-attention over the point cloud with Sinkhorn assignment \citep{sarlin2020superglue}, trained on synthetic data. Finger-marker labelling with spatial models and tracking \citep{alexanderson2017labeling}.
  \item \textbf{Cleaning and gap filling.} Learned denoising of optical data \citep{holden2018robust}; trajectory repair with a kinematic reference \citep{perepichka2019robust}; Vicon Nexus pattern, spline and rigid-body filling \citep{vicon_nexus}.
  \item \textbf{Gap in the literature.} No validated pipeline targets sparse hand markers during object interaction with participant-specific labelling and filling, and with a quantified account of each stage. \tbd{confirm with a literature search}
\end{itemize}

% ---------------------------------------------------------------------------------------------------------------
\section{Methods}

\subsection{Participants, protocol and marker set}
17 participants were recorded at 200\,Hz with Vicon \tbd{system and camera count}. The marker set has 22 markers:
4 on the forearm, 3 on the palm, 3 per finger and 3 on the thumb. Movement trials were \emph{Hands only} and
\emph{HOI}, in two sessions each, about 40--65k frames per trial \tbd{task description and objects}. Each participant
also performed a static trial. P1 and P2 used a reduced marker configuration and are handled separately. Nine
participants have accepted final statics (P3, P4, P6, P11, P12, P14, P15, P16, P17); for P17 only Trial 1 is
labelled. \tbd{ethics approval number; consent}

\subsection{Stage 1: personalised static MANO fit}
MANO \citep{romero2017mano} shape, global scale (about the wrist), pose and the placement of each marker on the mesh
surface are fitted to the static trial. Each marker attaches to a dorsal surface vertex of its segment and is held
3\,mm out along the vertex normal. The loss penalises marker-to-vertex distance above a cap, markers inside the skin,
and joints outside anatomical limits (finger twist, PIP/DIP abduction, DIP flexion coupled to PIP). Every fit is
reviewed visually by the first author and accepted or refitted. \tbd{review criteria and how many fits needed refits}

\paragraph{Photo-guided attachment.} Each marker's starting attachment (``pin'') is a MANO vertex; MANO's topology is
shared, so one pin set applies to every participant. The pins were set on one participant (P4): the first author
picked them by hand on the fitted mesh, and a photograph of the marked hand (dorsal view, same session as the static)
refined the palm and knuckle pins. The marker dots in the photo are detected, the camera is fitted to the near-rigid
markers (Palm1--3, Thumb1, Index1--Pinky1) with the focal length from the photo's metadata, and a ray from the camera
through each dot is cast onto the mesh. For participants whose photo could be matched, the photo's palm width at the
knuckle and wrist rows and the positions of the three finger webs were also used as fit targets for the palm. All
other participants start from P4's pins and let each pin move to a nearby dorsal vertex of the same segment (within
15\,mm). This adds manual effort per participant (photo check, pin review), counted in Table~\ref{tab:bench-cost}.
\tbd{which participants had a usable photo and which used photo targets; photo-to-mesh agreement of the accepted
statics: marker positions, palm widths, web positions and finger lengths, in mm}

\subsection{Stage 2: learned labelling}
Following SOMA \citep{ghorbani2021soma}, a transformer encoder over the unordered per-frame cloud scores each point
against the 22 labels, and log-space Sinkhorn normalisation with a learnable dustbin \citep{sarlin2020superglue}
gives a soft assignment that allows ghosts and missing labels. Training clouds are synthetic: a marker library is
calibrated on the participant's static by an earlier automatic calibration (marker attachments, residual offsets,
shape), not by the reviewed final static of Stage~1 (see Section~\ref{sec:personalisation}), and posed with real
library poses, with finger-block mixing, shape augmentation, marker drop-out, ghost points and random rotation. Whole trials
are labelled per frame, then a piece vote along continuous trajectories assigns one label per piece.
\tbd{model size, training steps, architecture table}

\subsection{Stage 3: learned gap filling}
\paragraph{Pre-processing.} A labelled position that jumps more than 8\,mm from its previous one is treated as a
label switch and dropped. Gaps of up to 10 frames are filled by a cubic spline.
\paragraph{Network.} Each frame is aligned to a template by a weighted Kabsch fit \citep{kabsch1976} on the near-rigid
markers (Palm1--3, Index1, Middle1, Ring1, Pinky1), weighted by visibility. The template is obtained by generalised
Procrustes analysis \citep{gower1975procrustes} on the training frames. A 4-layer GELU MLP (width 512) takes the
aligned coordinates of all markers and a visibility mask, and predicts every marker; the output is mapped back to
the laboratory frame. Frames with fewer than 3 visible near-rigid markers are not filled.
\paragraph{Training.} The filler is trained on the participant's own SOMA-labelled frames. In each training frame,
each observed marker is hidden with a probability drawn per frame from $U(0, 0.5)$, and the loss is the mean
Euclidean error on the hidden markers only. Incomplete frames are therefore usable, since each contributes its
observed markers. Training uses AdamW with a one-cycle schedule, 20{,}000 steps of 1{,}024 frames, seed 0.
\paragraph{Filling.} For a gap longer than 10 frames, the network predicts the marker in every gap frame from the
markers visible there; the prediction is shifted so that it meets the observed position at both gap edges, with the
shift blended linearly across the gap.

\subsection{Stage 4: movement fit}
For each frame, MANO pose, global orientation and translation are fitted with the participant's static shape, scale
and attachments held fixed. Each frame is initialised from the static pose moved rigidly onto that frame's palm
markers. The loss combines marker-to-attached-vertex distance (with the 3\,mm standoff), non-penetration, a cap on
marker-to-vertex distance, anatomical joint limits, and frame-to-frame smoothness. Each marker's terms are weighted
1 if observed and $w_\text{filled}=0.1$ if gap-filled. Frames are fitted in chunks of 1{,}000 with 800 Adam
iterations. In the runs reported here, every 10th frame of every trial was fitted.

\subsection{Stage 5: diagnostics and viewer}
Frames whose fitted pose leaves the anatomical limits or whose joint speed or rotation is implausible are flagged.
An interactive 3-D viewer shows markers, mesh and flags for every fit. \tbd{thresholds}

\subsection{Personalisation across stages}\label{sec:personalisation}
\paragraph{Design principle.} The participant's reviewed final static (Stage~1: shape, global scale and marker
attachments) is the single source of personalisation. It is meant to be used in every stage that needs a hand
model:
\begin{itemize}[nosep]
  \item the synthetic training clouds for the labeller (a library built from the final static);
  \item optionally, synthetic training clouds for the gap filler;
  \item the movement fit.
\end{itemize}
Whether personalisation improves each stage is treated as a hypothesis and tested
(Section~\ref{sec:does-personalisation-matter}), not assumed.
\paragraph{Current state.} The movement fit (Stage~4) uses the final static. The labeller (Stage~2) was trained on
a library from an earlier automatic static calibration. The gap filler (Stage~3) uses no hand model; it learns from
the participant's SOMA-labelled frames.

\subsection{Reproducibility}
The pipeline is scripted end to end. Training is seeded; two repeated trainings gave bit-identical weights on the
same GPU. Each fitted trial writes a provenance record with the code commit, uncommitted changes, hashes of the code
that ran and hashes of every input (labels, final static, filler). The pin file and the photo-alignment scripts are
released with the code; the hand photographs are shared only if the consent and ethics approval allow it, and
otherwise only the derived pins and measurements. Code: \tbd{repository URL, licence, DOI; photo consent}.

% ---------------------------------------------------------------------------------------------------------------
\section{Validation protocol}
\begin{enumerate}[nosep]
  \item \textbf{Ground truth.}
  \begin{itemize}[nosep]
    \item 500 hand-verified frames, stratified across participants and weighted toward hard cases.
    \item One fully verified trial, P10 Trial2 Hands only (about 42k frames). Its gap-fills are unreliable; 94\% of its values match raw Vicon points, and only those are used.
    \item \tbd{Continuous hand-verified HOI grasp windows of 5--10\,s, 2--3 per participant, at least 5 participants.}
  \end{itemize}
  \item \textbf{Labelling.} SOMA vs Vicon labels on the 500 frames and the verified trial; per participant, per marker, hands-only vs HOI. Participants whose data trained a model are held out of its evaluation.
  \item \textbf{Gap filling, artificial gaps.} Gaps of 1, 5, 10 and 25\,s are cut where the marker was observed (8 markers, up to 6 cuts per marker and length, fixed seed), refilled and compared with the true positions (median and worst error). Fillers trained on a participant's data are evaluated leave-one-trial-out.
  \item \textbf{Gap filling, real occlusions.} \tbd{inside the verified grasp windows: own-data filler vs rigid fill vs synthetic filler vs Vicon/Nexus filling}
  \item \textbf{Completeness.} Fraction of missing values each method leaves empty.
  \item \textbf{Movement fit.} Fit error on observed markers only (markers whose value is a real observation), so that methods with different fills are compared on the same markers: per frame the largest and the median marker-to-surface distance, median over frames; plus flagged-frame rate.
  \item \textbf{Sensitivity.} $w_\text{filled} \in \{0.03, 0.1, 0.3\}$ \tbd{}.
  \item \textbf{Optional angle reference.} \tbd{goniometer, video or extra markers, if the venue requires angle accuracy}
  \item \textbf{External benchmark on dense-marker data.} \tbd{} Take a public hand--object data set recorded with
  dense markers and reference MANO fits, for example GRAB \citep{taheri2020grab}. Subsample its markers to the
  22-marker layout and run the pipeline. Compare the recovered joint angles and gap fills with the data set's reference
  fits; this gives an angle reference without new recordings. To check: \tbd{licence}, \tbd{marker-layout mapping
  to the 22-marker set}, and \tbd{whether the data set's sessions resemble our task battery}.
\end{enumerate}

\subsection{Does personalisation matter?}\label{sec:does-personalisation-matter}
\tbd{all three tests}
\begin{enumerate}[label=(\alph*),nosep]
  \item \textbf{Synthetic test with known ground truth.} Generate marker trajectories from known hand shapes and
  attachments. Fit them with (i) the matched shape, (ii) the MANO mean shape (betas 0, default scale) and (iii) a
  mismatched participant's shape. Report the joint-angle error per joint, separately for proximal and distal joints.
  \item \textbf{Real-data test.} Fit the same trials (at least 3 participants, Hands only and HOI) with three
  statics: the MANO mean shape, the earlier automatic calibration and the reviewed final static. Report:
  \begin{itemize}[nosep]
    \item fit error on observed markers;
    \item the difference in joint angles between the statics;
    \item flagged-frame rates.
  \end{itemize}
  \item \textbf{Labeller personalisation.} Compare labelling accuracy on the 500 verified frames (only the reviewed
  frames), per participant, starting with P16, for a generic labeller, the labeller from the earlier static and the
  labeller from the final static. Repeat each training under identical seeded settings to estimate run-to-run
  variance. Also test a fix for the possible over-stretching (Section~\ref{sec:res-personalisation}), for example
  scale-only personalisation or limited shape coefficients for synthetic training, and keep the segment-length check as
  a script.
\end{enumerate}

\subsection{Benchmarking}\label{sec:benchmarking}
Each stage is compared with the alternatives a laboratory would realistically use, on the same data and against the
same ground truth (Table~\ref{tab:benchplan}). \tbd{all benchmarks}

\begin{table}[h]\centering\small
\caption{Benchmarking plan: one comparison per pipeline stage, on the same data.}\label{tab:benchplan}
\begin{tabularx}{\linewidth}{@{}p{2.1cm}XX@{}}\toprule
Stage & Compared methods & Metric \\\midrule
Labelling & Vicon's raw labels; a generic SOMA labeller (P10's model); the participant-calibrated SOMA labeller;
optionally the repository's other labellers (spatial GMM with Viterbi \citep{alexanderson2017labeling}, 3-D U-Net) &
Accuracy on the verified frames (reviewed frames only), Hands only vs HOI, per participant \\
Gap filling & Vicon/Nexus spline, pattern and rigid-body filling \citep{vicon_nexus}; the rigid-body fill implemented
here; the synthetic-trained filler; the own-data filler & Error inside verified HOI grasp windows; error on artificial
gaps, all methods on identical cuts; share of missing values left unfilled \\
Fitting & Rigid vs learned fill; filled values unweighted vs weighted; MANO mean shape vs earlier automatic
calibration vs reviewed final static (Section~\ref{sec:does-personalisation-matter}) & Fit error on observed markers;
flagged-frame rate; joint-angle differences \\
Practical cost & Every stage & Runtime and GPU memory per trial; training time per participant; manual review time
per participant (static review, label checks) \\\bottomrule
\end{tabularx}
\end{table}

\paragraph{Out of scope.} A full MoSh++ \citep{loper2014mosh,mahmood2019amass} or SOMA \citep{ghorbani2021soma}
reimplementation on sparse markers is not attempted. Both are heavyweight and bound by licences, and they target
dense body-and-hand marker sets, so the comparison would not be like-for-like. They are cited, and the differences in
marker density, personalisation and target task are explained in the Related Work section.

% ---------------------------------------------------------------------------------------------------------------
\section{Results}
\note{Tables~\ref{tab:gapfill} and~\ref{tab:fit} are from one participant (P16); the cohort results for all nine participants follow them.}

\subsection{Label quality of the raw data}
Across 17 participants, 0.1--15\% of frames per trial were fully geometrically self-consistent under Vicon's labels;
HOI trials were 2--10$\times$ worse than hands-only trials.

\subsection{Labelling accuracy}
\tbd{SOMA vs Vicon on the 500 verified frames and on P10 Trial2 Hands only; per participant, per marker, hands-only vs HOI (Table~\ref{tab:labelling}).}

\subsection{Personalising the labeller: no improvement on the monitor, cause under investigation}\label{sec:res-personalisation}
P16's labeller was retrained on a library built from the reviewed final static instead of the earlier automatic
calibration. It did not improve on the training monitor: Sinkhorn-assignment accuracy went from 85.0\% to 80.5\%,
and plain argmax accuracy from 87.9\% to 86.1\%. More points were left unlabelled. This is not evidence that
personalising the labeller hurts, for five reasons:
\begin{itemize}[nosep]
  \item \textbf{The monitor truth is mostly unverified.} The monitor uses 2000 frames with at most 4 missing markers,
  drawn at random from P16's Trial 1 Hands-only manually labelled export. That export is trustworthy only at the
  frames reviewed for the 500-frame sample, of which only 3 lie in this trial; just 1 of the 2000 monitor frames was
  reviewed. Elsewhere the export carries Vicon-like labels, so the monitor mostly measures agreement with those.
  \tbd{scores of both labellers on P16's reviewed frames}
  \item \textbf{One run each.} Both models used the same training code, step count (10{,}000), seed and monitor
  frames, and differed only in the library. Run-to-run variance is unknown \tbd{repeat seeded trainings}.
  \item \textbf{The fit results are mixed.} On observed markers, the largest error per frame was lower with the
  retrained labels on both Hands-only trials and Trial2 HOI, and higher on Trial 1 HOI; the typical error was lower on
  both Hands-only trials and higher on both HOI trials. The two label sets observe different markers, so this is
  not a like-for-like comparison.
  \item \textbf{More unlabelled points are ambiguous.} They may reflect more caution about ghost points rather than
  more errors; this cannot be decided without verified truth.
  \item \textbf{The proposed cause is untested.} The hypothesis below rests on ad-hoc segment-length checks that were
  not kept as scripts.
\end{itemize}
\paragraph{Hypothesis.} \tbd{to be tested in Section~\ref{sec:does-personalisation-matter}(c)} The final static's
large shape coefficients, posed with another participant's library poses, produce over-stretched synthetic hands.
The final static would then be a good fit to the static pose, but a poor generator of training poses.
\tbd{results of tests (a)--(c)}

\begin{table}[h]\centering\small
\caption{Labelling accuracy (\% of points correctly labelled). \tbd{all values}}\label{tab:labelling}
\begin{tabular}{lcccc}\toprule
 & \multicolumn{2}{c}{500 verified frames} & \multicolumn{2}{c}{P10 verified trial} \\
 & Vicon & SOMA & Vicon & SOMA \\\midrule
Hands only & \tbd{} & \tbd{} & \tbd{} & \tbd{} \\
HOI & \tbd{} & \tbd{} & -- & -- \\\bottomrule
\end{tabular}
\end{table}

\subsection{Gap-fill accuracy}
\begin{table}[h]\centering\small
\caption{Gap-fill error on artificial gaps cut where the marker was observed, P16, median / worst (mm).
Own-data fillers are trained on the participant's other trials (test trial held out).}\label{tab:gapfill}
\begin{tabular}{llccc}\toprule
Gap & Trial & Rigid-body fill & Synthetic-trained filler & Own-data filler \\\midrule
1\,s & Trial2 HOI & -- & 3.6 / 26 & 2.4 / 10 \\
10\,s & Trial2 HOI & -- & 2.5 / 7.6 & 1.5 / 2.4 \\
25\,s & Trial 1 HOI & 3.8 / 11.6 & 5.6 / 8.3 & 3.5 / 4.5 \\\bottomrule
\end{tabular}
\end{table}
\note{The 25\,s rigid value comes from a different set of artificial cuts than the two filler values. Before submission, rerun all methods on identical cuts, fill the rigid column for every row, and add the Vicon/Nexus baseline on the same cuts (Table~\ref{tab:bench-fill}).}

At 5\,s and 10\,s gaps, where the rigid fill was run, median errors were 2.4--14.0\,mm for the rigid fill and
1.8--5.9\,mm for the synthetic-trained filler. The own-data filler's median was below the synthetic filler's on
Trial2 HOI, but not on Trial 1 HOI. Worst errors were 20--32\,mm with the synthetic filler and 7--10\,mm with the
own-data filler. A filler trained on another participant (P10) transferred poorly, worse than the synthetic filler.

\paragraph{Cohort.} For each of the other eight participants, a filler was trained without one HOI trial and tested on
it against the rigid fill, on identical cuts (Table~\ref{tab:gapfill-cohort}).
\begin{itemize}[nosep]
  \item The own-data filler lowered the worst error for seven of eight participants, for example from 44.9 to 9.8\,mm
  at 10\,s for P12.
  \item Its median at gaps of 5\,s and longer was lower for every participant except P17.
  \item It left no cut empty, where the rigid fill left 0--4 per gap length.
  \item P17, whose filler was trained on a single other trial, was the exception: the two methods were about equal.
\end{itemize}
\begin{table}[h]\centering\small
\caption{Held-out gap-fill error, median / worst (mm), rigid fill $\rightarrow$ own-data filler, identical cuts in each
participant's held-out HOI trial (Trial2 HOI; Trial 1 HOI for P17).}\label{tab:gapfill-cohort}
\begin{tabular}{lcccc}\toprule
P & 1\,s & 5\,s & 10\,s & 25\,s \\\midrule
P3 & 2.9/32.0 $\rightarrow$ 1.9/9.7 & 3.9/17.4 $\rightarrow$ 1.8/6.3 & 4.0/34.9 $\rightarrow$ 3.0/9.9 & 8.2/25.0 $\rightarrow$ 3.6/6.2 \\
P4 & 1.8/25.5 $\rightarrow$ 1.8/17.4 & 7.5/32.6 $\rightarrow$ 3.4/16.5 & 11.2/32.1 $\rightarrow$ 4.1/16.7 & 14.7/22.9 $\rightarrow$ 5.9/14.9 \\
P6 & 2.1/67.2 $\rightarrow$ 1.7/28.8 & 4.4/39.2 $\rightarrow$ 3.3/29.4 & 5.4/32.3 $\rightarrow$ 4.2/28.7 & -- \\
P11 & 3.5/28.7 $\rightarrow$ 2.3/13.4 & 4.7/27.4 $\rightarrow$ 3.1/13.6 & 6.2/22.0 $\rightarrow$ 5.2/13.3 & 11.8/21.0 $\rightarrow$ 3.8/11.7 \\
P12 & 2.9/33.0 $\rightarrow$ 1.9/17.2 & 5.3/53.1 $\rightarrow$ 2.8/14.0 & 9.9/44.9 $\rightarrow$ 4.5/9.8 & 11.4/24.0 $\rightarrow$ 5.1/5.4 \\
P14 & 1.5/44.3 $\rightarrow$ 1.9/9.8 & 2.8/42.9 $\rightarrow$ 2.0/11.0 & 12.7/40.1 $\rightarrow$ 5.0/10.5 & 5.6/24.2 $\rightarrow$ 4.3/6.9 \\
P15 & 4.4/25.7 $\rightarrow$ 2.8/20.1 & 9.1/24.8 $\rightarrow$ 3.6/17.0 & 8.4/47.3 $\rightarrow$ 3.6/14.9 & 8.0/47.3 $\rightarrow$ 4.8/13.5 \\
P17 & 1.5/21.2 $\rightarrow$ 1.8/24.5 & 2.2/21.8 $\rightarrow$ 2.1/24.6 & 3.5/22.3 $\rightarrow$ 3.7/23.2 & 3.0/21.3 $\rightarrow$ 2.6/23.2 \\\bottomrule
\end{tabular}
\end{table}
\note{Cohort values measured 2026-10-29 (logs: results/soma/real\_fill/<P>\_heldout.log; P3's rigid row from the rerun after the empty-cut fix).}
\tbd{errors inside verified grasp windows; Vicon/Nexus filling as a baseline}

\subsection{Completeness}
The rigid fill left about 50--60\% of missing values empty. The learned filler left 0.7--27\% empty per trial; what
remains is frames in which the hand is not visible. Across all 34 trials of the nine participants, the rigid fill
left 1.6--76.4\% of missing values empty (median 47.6\%) and the learned filler 0.0--39.1\% (median 10.2\%). The
learned filler left fewer values empty in 29 trials, the same in 1, and more in 4 Hands-only trials (by 1.1--8.5
points), probably frames with fewer than 3 near-rigid markers visible, which it does not fill \tbd{check}.

\subsection{Movement fit}
\begin{table}[h]\centering\small
\caption{Fit error on observed markers only, P16, median over frames (mm): rigid fill $\rightarrow$ learned fill with filled values weighted 0.1. Two trials per condition.}\label{tab:fit}
\begin{tabular}{lcc}\toprule
Condition & Largest per frame & Typical per frame \\\midrule
Hands only & 12.2 $\rightarrow$ 11.4;\; 11.7 $\rightarrow$ 11.5 & 6.4 $\rightarrow$ 6.3;\; 6.7 $\rightarrow$ 6.5 \\
HOI & 14.3 $\rightarrow$ 14.7;\; 14.7 $\rightarrow$ 15.2 & 7.6 $\rightarrow$ 7.5;\; 7.8 $\rightarrow$ 7.7 \\\bottomrule
\end{tabular}
\end{table}
With the filled values unweighted, HOI fits were about 1\,mm worse on observed markers; down-weighting closed most of
that gap.

\paragraph{Cohort.} All nine participants were fitted in full (every 10th frame of every trial, 167 batches, none
failed), with the own-data filler and $w_\text{filled}=0.1$, and compared with the rigid-fill fits on the same batches
(Table~\ref{tab:fit-cohort}).
\begin{itemize}[nosep]
  \item Hands only: the learned fill was never worse than the rigid fill.
  \item HOI: better for five participants (P3, P4, P11, P12, P14), and worse on the largest error for four (P6, P15,
  P16, P17), by up to 0.7\,mm.
\end{itemize}
\begin{table}[h]\centering\small
\caption{Change in observed-marker fit error, learned (weighted 0.1) minus rigid fill, nine participants, 17 trials per
condition: median change in mm (range), and trials better / same / worse.}\label{tab:fit-cohort}
\begin{tabular}{lcccc}\toprule
 & \multicolumn{2}{c}{Largest per frame} & \multicolumn{2}{c}{Typical per frame} \\
Condition & Change & B / S / W & Change & B / S / W \\\midrule
Hands only & $-0.4$ ($-0.9$ to $0.0$) & 14 / 3 / 0 & $-0.1$ ($-0.2$ to $0.0$) & 9 / 8 / 0 \\
HOI & $-0.4$ ($-1.6$ to $+0.7$) & 10 / 0 / 7 & $-0.2$ ($-0.3$ to $+0.4$) & 13 / 2 / 2 \\\bottomrule
\end{tabular}
\end{table}
\note{Cohort values measured 2026-10-29 (logs: results/soma/real\_fill/<P>\_score.log).}
\tbd{flagged-frame rates; weight sweep}
\note{Table~\ref{tab:fit} will be extended with fits on Vicon/Nexus-filled data (Table~\ref{tab:bench-fit}).}

\subsection{Benchmarks}\label{sec:res-bench}
\tbd{all values; method and plan in Section~\ref{sec:benchmarking}}

\begin{table}[h]\centering\small
\caption{Labelling benchmark on the nine participants with a reviewed final static (P3, P4, P6, P11, P12, P14--P17):
share of real marker points named correctly on the hand-reviewed frames (\%). First number: all points pooled; in
brackets: median over participants. Hands only: 72 frames, 1462 points; HOI: 87 frames, 1695 points.}\label{tab:bench-label}
\begin{tabular}{lcc}\toprule
Method & Hands only & HOI \\\midrule
Vicon raw labels & 92.9 (95.1) & 80.9 (86.4) \\
Generic SOMA (P10 model) & 97.1 (96.2) & 86.4 (92.3) \\
Participant-calibrated SOMA & 98.2 (98.7) & 90.1 (93.4) \\
Spatial GMM + Viterbi (optional) & \tbd{} & \tbd{} \\
3-D U-Net (optional) & \tbd{} & \tbd{} \\\bottomrule
\end{tabular}
\end{table}
The calibrated labeller was the most accurate method in both conditions: 98.2\% against 92.9\% for Vicon on Hands
only, and 90.1\% against 80.9\% on HOI, with the generic labeller in between. Per participant, it beat Vicon for 7 of 9
participants on Hands only and 6 of 9 on HOI, and beat the generic labeller for 5 and 7. Vicon was better on HOI for
P14, P15 and P16 (P16's mirrored-hand episode, Section~\ref{sec:res-personalisation}), and on Hands only for P6 and
P11. P6 Trial 1 HOI was not scored, because its corrected export has a different number of frames from the raw file.
Across all 17 reviewed participants the pattern is the same on HOI; on Hands only, one trial outside this set (P9 Trial2
Hands only, 63 frames) scores 52\% for both SOMA labellers against 92\% for Vicon, which is not yet investigated.
\note{Measured 2026-10-29 with scripts/soma/eval\_sample.py (results/soma/eval/bench\_\{generic,calibrated\}\_*). SOMA
scores are after the piece vote. Spatial GMM and 3-D U-Net not yet run.}

\begin{table}[h]\centering\small
\caption{Gap-filling benchmark: median / worst error (mm), all methods on identical cuts, and share of missing values
left unfilled. \tbd{}}\label{tab:bench-fill}
\begin{tabular}{lccc}\toprule
Method & HOI grasp windows & Artificial gaps & Left unfilled \\\midrule
Vicon/Nexus spline & \tbd{} & \tbd{} & \tbd{} \\
Vicon/Nexus pattern & \tbd{} & \tbd{} & \tbd{} \\
Vicon/Nexus rigid body & \tbd{} & \tbd{} & \tbd{} \\
Rigid-body fill (ours) & \tbd{} & \tbd{} & \tbd{} \\
Synthetic-trained filler & \tbd{} & \tbd{} & \tbd{} \\
Own-data filler & \tbd{} & \tbd{} & \tbd{} \\\bottomrule
\end{tabular}
\end{table}

\begin{table}[h]\centering\small
\caption{Fitting benchmark: fit error on observed markers (mm), flagged-frame rate (\%) and joint-angle difference
(deg) relative to \tbd{reference configuration}. \tbd{}}\label{tab:bench-fit}
\begin{tabular}{lccc}\toprule
Configuration & Fit error & Flagged frames & Angle difference \\\midrule
Rigid fill & \tbd{} & \tbd{} & \tbd{} \\
Learned fill, unweighted & \tbd{} & \tbd{} & \tbd{} \\
Learned fill, weighted 0.1 & \tbd{} & \tbd{} & \tbd{} \\
Vicon/Nexus fill & \tbd{} & \tbd{} & \tbd{} \\
MANO mean shape & \tbd{} & \tbd{} & \tbd{} \\
Earlier automatic calibration & \tbd{} & \tbd{} & \tbd{} \\
Reviewed final static & \tbd{} & \tbd{} & \tbd{} \\\bottomrule
\end{tabular}
\end{table}

\begin{table}[h]\centering\small
\caption{Practical cost per stage. \tbd{}}\label{tab:bench-cost}
\begin{tabular}{lcccc}\toprule
Stage & \shortstack{Runtime\\per trial} & \shortstack{GPU\\memory} & \shortstack{Training per\\participant} & \shortstack{Manual review\\per participant} \\\midrule
Static fit (incl.\ photo alignment, pin review) & \tbd{} & \tbd{} & -- & \tbd{} \\
Labelling & \tbd{} & \tbd{} & \tbd{} & \tbd{} \\
Gap filling & \tbd{} & \tbd{} & \tbd{} & -- \\
Movement fit & \tbd{} & \tbd{} & -- & -- \\
Diagnostics and viewer & \tbd{} & \tbd{} & -- & \tbd{} \\\bottomrule
\end{tabular}
\end{table}

\subsection{External dense-marker benchmark}\label{sec:res-external}
\tbd{joint-angle and gap-fill error against the reference fits of a public dense-marker hand--object data set,
subsampled to the 22-marker layout}

% ---------------------------------------------------------------------------------------------------------------
\section{Discussion}
The benchmarks quantify where each stage beats the standard alternative and where it does not; that per-stage
account, not a single headline number, is the paper's contribution.
\begin{itemize}[nosep]
  \item \textbf{Where each stage helps.} Learned labelling \tbd{}; participant-specific filling mainly cuts worst-case errors and completes trials; down-weighting makes complete trials compatible with fit quality on real markers.
  \item \textbf{Where a stage does not help.} A filler from another participant; unweighted filled values in HOI.
  \item \textbf{Inconclusive.} Retraining the labeller on the final static did not improve on the monitor; the cause is under investigation (Section~\ref{sec:res-personalisation}).
  \item \textbf{Why own-data training wins.} Marker placement and hand geometry are participant-specific; a filler trained on another participant transfers that participant's offsets.
  \item \textbf{Personalisation from the first stage.} If personalisation helps, it should be applied from the first stage, so that every stage sees the same hand. \tbd{depends on Section~\ref{sec:does-personalisation-matter}}
  \item \textbf{Open question.} Personalisation that is good for fitting might be harmful for synthetic data generation, if it produces unrealistic training clouds; the labeller result (Section~\ref{sec:res-personalisation}) is consistent with this but does not show it.
  \item \textbf{Practical recommendation.} \tbd{once cohort and grasp-window results are in}
\end{itemize}
\paragraph{Limitations.}
\begin{itemize}[nosep]
  \item Occluded markers have no direct ground truth; artificial gaps are cut where markers are visible, mostly in easier poses, so they are optimistic for real grasp occlusions. The planned grasp windows address this.
  \item Label quality bounds filler quality: the filler learns from the participant's SOMA labels, including their errors.
  \item One filled-value weight (0.1) has been tested, on one participant.
  \item Gap-fill accuracy is measured on artificial gaps only; on HOI the fit with learned filling is still worse than with the rigid fill for four of nine participants.
  \item Fit error measures agreement with markers, not joint-angle accuracy; no independent angle reference \tbd{}.
  \item Without an independent angle reference, the real-data personalisation test shows that the angles change with the static, not which angles are correct. Only the synthetic test measures correctness.
  \item The personalised static relies on manual review.
\end{itemize}

\section{Conclusion}
\tbd{two to three sentences once results are final}

% ---------------------------------------------------------------------------------------------------------------
\section*{Figures and tables}
\begin{tabularx}{\linewidth}{@{}lX@{}}\toprule
Item & Content \\\midrule
Fig.~1 & Pipeline overview: static fit $\rightarrow$ labelling $\rightarrow$ gap filling $\rightarrow$ movement fit $\rightarrow$ diagnostics, with inputs, outputs and provenance. \\
Fig.~2 & Marker set on the hand and the personalised MANO static fit (one participant), attachments shown. \\
Fig.~3 & Raw label quality: self-consistent-frame fraction per trial, hands-only vs HOI, all participants. \\
Fig.~4 & Labelling examples: a Vicon mislabel and the SOMA relabel (crossed vs straight finger chains). \\
Fig.~5 & Gap filling: error vs gap length for rigid, synthetic-trained and own-data fillers (median and worst). \\
Fig.~6 & A long HOI gap: true, rigid-filled and learned-filled trajectory of one fingertip marker. \\
Fig.~7 & Movement-fit error on observed markers per participant and condition, rigid vs learned fill (cohort). \\
Fig.~8 & Filled-value weight sweep: observed-marker error vs completeness. \\
Fig.~9 & Viewer screenshot with flagged frames. \\
Fig.~10 & Joint-angle error vs shape mismatch (synthetic), and between-static angle differences on real trials. \tbd{} \\
Fig.~11 & Benchmark summary: each stage against its baselines (labelling accuracy, fill error, fit error). \tbd{} \\
Fig.~12 & Practical cost per stage (runtime, GPU, manual review time). \tbd{} \\
Fig.~13 & External benchmark: joint-angle error on subsampled dense-marker data. \tbd{} \\
Tab.~1 & Participants, trials, frames, marker availability. \\
Tab.~2 & Labelling accuracy (Table~\ref{tab:labelling}). \\
Tab.~3 & Gap-fill accuracy (Table~\ref{tab:gapfill}), extended to the cohort and grasp windows. \\
Tab.~4 & Completeness per method. \\
Tab.~5 & Movement-fit error (Table~\ref{tab:fit}), extended to the cohort. \\
Tab.~6 & Negative and inconclusive results: cross-participant filler; labeller retrained on the final static (monitor, 500-frame scores, repeat runs). \\\bottomrule
\end{tabularx}

% ---------------------------------------------------------------------------------------------------------------
\section*{Expected reviewer objections}
\begin{tabularx}{\linewidth}{@{}XX@{}}\toprule
Objection & Answer / planned validation \\\midrule
``No ground truth for occluded markers; artificial gaps are easy cases.'' & Verified HOI grasp windows (5--10\,s, 2--3 per participant, $\geq$5 participants) evaluate filling on real occlusions. \\
``One participant.'' & Cohort run on all nine participants with accepted statics. \\
``Just SOMA/MoSh for hands.'' & The contribution is the validated combination for sparse markers during HOI; stage-by-stage results, including negative ones, show what each part adds. \\
``Why not just use Vicon/Nexus gap filling?'' & Nexus spline, pattern and rigid-body filling are baselines in the gap-filling benchmark, on identical cuts and in the grasp windows (Table~\ref{tab:bench-fill}). \tbd{} \\
``No joint-angle ground truth.'' & External benchmark on subsampled dense-marker data (Section~\ref{sec:res-external}), alongside the synthetic personalisation test (Section~\ref{sec:does-personalisation-matter}a). \tbd{} \\
``Is it practical for a small lab?'' & Practical-cost table: runtime, GPU memory, training and manual review time per stage (Table~\ref{tab:bench-cost}). \tbd{} \\
``The 0.1 weight is arbitrary.'' & Weight sweep (0.03, 0.1, 0.3). \\
``Fit error is not angle accuracy.'' & Optional independent angle reference; otherwise scope the claims to marker-level accuracy. \\
``The filler learns from SOMA labels; errors propagate.'' & Report labelling accuracy per participant and relate it to filler error. \\
``Manual static review is subjective.'' & Report review criteria and refit counts; release the static fits and views. \\
``Is personalisation necessary? Would a generic MANO do?'' & Tests (a)--(c) in Section~\ref{sec:does-personalisation-matter}: synthetic angle error vs shape mismatch, real fits with three statics, labellers with and without personalisation. \tbd{} \\
``Reproducibility.'' & Seeded training (bit-identical repeats), provenance records, open code and tutorial. \\\bottomrule
\end{tabularx}

% ---------------------------------------------------------------------------------------------------------------
\appendix
\section{Tutorial: running the pipeline on a new participant}
\paragraph{Requirements.}
Python environment from the repository (\texttt{pip install -e ".[dev,train]"}, plus \texttt{plotly} and
\texttt{matplotlib}); a CUDA GPU recommended. MANO weights must be obtained separately from the MANO website under
its licence; the scripts look for them in a sibling \texttt{clean\_kinematics/mano\_v1\_2/models} directory, else
under the data root, or take \texttt{--mano-dir}. Data sharing is subject to ethics approval \tbd{}. Trial files are
located through \texttt{manual\_frames.csv} (participant, session, CSV name).

\paragraph{1. Static fit (reviewed by hand).}
\begin{lstlisting}
python scripts/soma/static_search.py ...            # search fit settings for the participant  [TBD exact arguments]
python scripts/soma/static_final_fit.py --participant P --start <search result> --out OUT1
python scripts/soma/static_final_fit.py --participant P --start OUT1 --rounds 3 --out OUT2
# review OUT2/static.html; accept -> soma_transfer/static_final/P/P_static_final_params.npz
\end{lstlisting}

\paragraph{2. Labelling (calibrate, train, label every trial).}
\begin{lstlisting}
python scripts/soma/run_fullset.py --participants P     # -> results/soma/fullset/P/<trial>/soma_labelled.csv
\end{lstlisting}

\paragraph{3--5. Gap filling, movement fit, diagnostics and views.}
\begin{lstlisting}
bash scripts/soma/run_realfill.sh P
# held-out filler check  -> results/soma/real_fill/P_heldout.log
# filler on all trials   -> results/soma/real_fill/filler_P_all.pt
# fill + fit + views     -> results/soma/real_fill_w01/movement/P, OneDrive results/soma/P_realfill_w01
# score vs rigid fill    -> results/soma/real_fill/P_score.log
\end{lstlisting}
Each trial folder contains a \texttt{fit\_trials\_provenance.json} with the commit, code and input hashes and the arguments.
\note{Check the exact arguments of static\_search.py and run\_fullset.py before submission.}

\bibliographystyle{plainnat}
\bibliography{references}
\end{document}
